表~\ref{tab:information-candidate-support}将支持方向与两种计算并列。完整恢复证据要求候选覆盖后验，
而有限的\(D_{\mathrm{KL}}(q\Vert\pi)\)还要求候选不在零后验区域放置质量；两者控制的对象不同。
\begin{table*}[htbp]
  \centering\small
  \begin{tabularx}{\linewidth}{XXcccc}
    \toprule
    后验\(\pi\) & 候选\(q\) & \(q\ll\pi\) & \(\pi\ll q\) & \(D_{\mathrm{KL}}(q\Vert\pi)\) & 证据恢复比例\\
    \midrule
    \((1/2,1/2)\) & \((1,0)\) & 是 & 否 & \(\log2\) & \(1/2\)\\
    \((1,0)\) & \((1/2,1/2)\) & 否 & 是 & \(+\infty\) & \(1\)\\
    \((1/2,1/2)\) & \((1/2,1/2)\) & 是 & 是 & \(0\) & \(1\)\\
    \bottomrule
  \end{tabularx}
  \caption{候选支持对KL与证据恢复的影响。两点分布共用同一取值集合，证据满足
  \(0<p_{\bm\theta}(x)<\infty\)。恢复比例为\(\mathbb E_q[p_{\bm\theta}(x,Z)/q(Z)]/p_{\bm\theta}(x)=\pi(\{q>0\})\)，
  比值仅在\(q>0\)处评价。表中的有限KL是具体例子的结果；一般空间中\(q\ll\pi\)不能单独保证KL有限。}
  \label{tab:information-candidate-support}
\end{table*}
